\documentclass[12pt]{article}
\usepackage[T1]{fontenc}
\usepackage[utf8]{inputenc}
\usepackage{lmodern}
\usepackage{textcomp}
\usepackage{enumitem}
\usepackage{geometry}
\geometry{margin=1in}
\begin{document}
\begin{center}
Answer Key - Religious Studies - K-12 Level Assessment 10 \
\small (Answers are ordered exactly as in the question paper.)
\end{center}

\section*{Multiple Choice}
\begin{enumerate}[label=\arabic*.]
\item \textbf{Answer:} A
\\ \textit{Options:} A) Islamic identity; B) Islamic modesty; C) Islamic obedience; D) Islamic seclusion
\\ \textit{Note:} The foremost reason for wearing the hijab in the present day is to express and affirm one's Islamic identity, reflecting personal and cultural beliefs in a diverse world.
\item \textbf{Answer:} C
\\ \textit{Options:} A) Varro; B) Plato; C) Cicero; D) Augustus
\\ \textit{Note:} Cicero advocated for civic religion in his work The Laws by emphasizing the importance of religious practices in promoting moral values and social cohesion within the Roman state.
\item \textbf{Answer:} A
\\ \textit{Options:} A) 36,000; B) 12,000; C) 24,000; D) 7,000
\\ \textit{Note:} The correct answer is 36,000 because historical texts indicate that Chandanbala, a prominent figure in certain religious traditions, was known to have a significant number of female followers dedicated to renouncing worldly life under her guidance.
\item \textbf{Answer:} D
\\ \textit{Options:} A) Householders who supported renouncers; B) Orthodox interpreters of the scriptures; C) Women ascetics who lived in communities; D) Temple-dwelling renouncers
\\ \textit{Note:} The caityavasis in Jaina traditions are correctly identified as temple-dwelling renouncers because they are ascetics who live in and around temples, dedicating their lives to spiritual practices and adherence to Jain principles.
\item \textbf{Answer:} C
\\ \textit{Options:} A) Gods; B) Yoga; C) Devotion; D) Religion
\\ \textit{Note:} Bhakti is often translated as devotion because it refers to a personal, loving relationship and commitment to a deity or divine figure, emphasizing emotional attachment and worship in various religious traditions, particularly in Hinduism.
\end{enumerate}

\section*{True / False}
\begin{enumerate}[label=\arabic*.]
\setcounter{enumi}{5}
\item \textbf{Answer:} False
\\ \textit{Options:} A) True; B) False
\\ \textit{Note:} Guru Ram Das was not responsible for the establishment of the Darbar Sahib in Amritsar; it was his successor, Guru Arjan Dev, who completed its construction.
\item \textbf{Answer:} False
\\ \textit{Options:} A) True; B) False
\\ \textit{Note:} The statement is false because Akhenaten is primarily known for promoting the worship of the sun disk Aten and diminishing the significance of other gods, but he did not specifically focus on diminishing the roles of Shu and Set in his religious reforms.
\item \textbf{Answer:} True
\\ \textit{Options:} A) True; B) False
\\ \textit{Note:} The Namokar Mantra is considered the most important prayer in Jainism as it expresses reverence for the five supreme beings and is central to the spiritual practice of Jains, emphasizing their core beliefs and values.
\item \textbf{Answer:} False
\\ \textit{Options:} A) True; B) False
\\ \textit{Note:} The statement is false because while Pachomius was an important figure in early Christian monasticism, he is not considered the first significant anchorite hermit, as figures like Anthony the Great predate him and are recognized as pioneers of the hermitic lifestyle.
\item \textbf{Answer:} True
\\ \textit{Options:} A) True; B) False
\\ \textit{Note:} In Buddhism, the term "Bhikshuni" specifically refers to a fully ordained female monastic, paralleling the male equivalent known as "Bhikshu," thus making the statement true.
\end{enumerate}

\section*{Long Form Response}
\begin{enumerate}[label=\arabic*.]
\setcounter{enumi}{10}
\item \textbf{Answer:} Ziran, a key concept in Zhuangzi's philosophy, refers to the natural state of being or spontaneity, where things exist and operate in their own way without forced interference. In this context, spontaneity means allowing events and actions to unfold naturally, in harmony with the Dao, or the way of the universe. This understanding encourages individuals to embrace their true nature and the inherent rhythms of life, promoting a sense of freedom and authenticity. By recognizing ziran, Daoism teaches that one should align with nature rather than resist it, leading to a deeper appreciation of existence and the interconnectedness of all things. Ultimately, spontaneity is a guiding principle that fosters peace and balance in the world.
\\ \textit{Note:} correct because it accurately defines ziran as spontaneity in Zhuangzi's philosophy, emphasizing the importance of allowing natural processes to unfold in harmony with the Dao, which fosters authenticity and a deeper connection to the rhythms of life in Daoism.
\item \textbf{Answer:} Female poets such as Ghosa, Apala, and Lopamurda made significant contributions to the early Vedas, highlighting the active role of women in Vedic society. Their hymns and verses reveal that women were not only participants but also respected figures in the spiritual and cultural life of their communities. For example, Ghosa is known for her devotion and expression of longing for the divine, while Apala's poetry reflects personal strength and individuality. These works demonstrate that women had a voice in religious discourse, challenging the notion that Vedic society was solely patriarchal. Overall, the contributions of these female poets illustrate a more complex and inclusive view of women's roles during this period.
\\ \textit{Note:} correct because it emphasizes the significant contributions of female poets in the early Vedas, illustrating how their works reflect women's active and respected roles in the spiritual and cultural aspects of Vedic society.
\end{enumerate}

\vfill
\noindent\textit{End of Answer Key}
\end{document}